La actividad propuesta en equipo fue sencilla, no se necesito de comandos muy rebuscados
para la creación de los Issues, realmente fue algo muy sencillo la creación del Issue
dentro de la terminal, solo se ocuparon los comandos propuestos dentro del documento y
después de la creación del Issue al insertar un código mal estructurado, y luego comenzar
a corregirlo, luego de arreglar ese código, solo se envié con un git add, con su respectivo
git commit y finalmente el git push para registrarlo dentro del repositorio.
La creación del Issue en la pagina web fue un proceso principalmente más tardado que
la creación del Issue en la terminal, porque también se siguieron los comandos planteados
en el documento de la practica. Las dificultades presentadas fue la visualización de las
etiquetas dentro de la terminal y aparte de errores comunes de escritura al realizar las
practicas dentro de la misma terminal, en conclusión, tanto en la pagina web como en la
terminal, para agregar las etiquetas no se entienden muy bien esas partes, y como deben
de agregarse.